%=====================================================================
%  Appendix C.x : Primitive sufficient conditions for threshold stability
%
%  Drop-in subsection for the Supplementary Material of
%  "Factor-Adjusted Knockoff Inference for High-Dimensional Time Series"
%
%  Insert immediately after Assumption C.3.  Renumber C.x / PC / Lemma /
%  Proposition labels to fit the surrounding appendix.
%
%  Notation follows the main text:
%    \mathcal{G}   conditioned-on design sigma-field
%    G             the knockoff+ threshold of (3.18)   [note the clash;
%                  consider renaming the threshold to \mathcal{T} globally]
%    \zeta         = T^{-1/2} U' eps^dagger, the 2kappa-dim score
%    \zeta^\star   | \mathcal{G} ~ N(0,\Xi), the covariance-matched law
%    D             = T^{-1/2} U' g, the missed-signal drift
%    \Delta_{NT}   coupling scale from Assumption C.2
%    W^\star       LCD statistics computed from D + \zeta^\star
%=====================================================================

\subsection{Primitive sufficient conditions for threshold stability}
\label{sec:threshold-primitive}

Assumption~C.3 is stated at the level of the realized null counts under the
reference law $P_{D^\star}$.  This subsection decomposes it into three
conditions with separate interpretations, proves two of them from structure
already imposed elsewhere in the paper, and isolates the single remaining
requirement.  The purpose is to make transparent exactly which part of
Assumption~C.3 is a genuine primitive restriction and which part is
bookkeeping.

Throughout, $\mathcal{H}_0 = \widetilde S \cap S_0^{c}$ denotes the retained
null set and $\kappa_0 = |\mathcal{H}_0|$.  Recall from Assumption~C.3 that,
under $P_{D^\star}$,
\[
V^{\star,+}(t) = \bigl|\{ j \in \mathcal{H}_0 : W_j^\star \ge t \}\bigr|,
\qquad
V^{\star,-}(t) = \bigl|\{ j \in \mathcal{H}_0 : W_j^\star \le -t \}\bigr|,
\]
\[
M^\star(t) = V^{\star,+}(t) + V^{\star,-}(t),
\qquad
\widetilde N^\star_{2\Delta}(t)
 = \bigl|\{ j \in \mathcal{H}_0 : \bigl| |W_j^\star| - t \bigr| \le 2\Delta_{NT} \}\bigr| .
\]
Define the conditional marginal null tail
\begin{equation}
\label{eq:Gbar}
\bar G(t) \;=\; \kappa_0^{-1} \sum_{j \in \mathcal{H}_0}
              P_{D^\star}\bigl( |W_j^\star| \ge t \,\big|\, \mathcal{G} \bigr),
\end{equation}
which is well defined because, conditional on $\mathcal{G}$, the pair-swap
invariance of $P_{D^\star}$ makes each null $W_j^\star$ symmetric about zero.
Write $I_{NT} = [\,v_{NT}\Delta_{NT}/2,\ \infty)$ for the effective threshold
range enlarged by one band width.

\subsubsection*{Three primitive conditions}

\begin{condition}[Relative anti-concentration of ideal null statistics]
\label{pc:1}
There is a sequence $\alpha_{NT} \to 0$ such that
\[
\sup_{t \in I_{NT}}\;
\frac{\kappa_0^{-1}\sum_{j \in \mathcal{H}_0}
      P_{D^\star}\bigl( \bigl||W_j^\star| - t\bigr| \le 2\Delta_{NT}
                        \,\big|\, \mathcal{G}\bigr)}
     {\bar G(t)}
\;\le\; \alpha_{NT}.
\]
\end{condition}

\begin{condition}[Conditional concentration of null counts]
\label{pc:2}
There are $m_{NT}$ with $m_{NT}/\ln(NT) \to \infty$ and a constant $C>0$ such
that, conditional on $\mathcal{G}$ and uniformly in $t \in I_{NT}$,
\[
\operatorname{Var}_{D^\star}\Bigl( \sum_{j \in \mathcal{H}_0}
     \mathbf 1\{ W_j^\star \ge t \} \,\Big|\, \mathcal{G} \Bigr)
\;\le\; C\, m_{NT}\, \kappa_0 \bar G(t)
   \;+\; o\bigl( \{\ln (NT)\}^{-1} [\kappa_0 \bar G(t)]^2 \bigr),
\]
and the same bound holds with $\{W_j^\star \ge t\}$ replaced by
$\{ \bigl||W_j^\star|-t\bigr| \le 2\Delta_{NT}\}$.
\end{condition}

\begin{condition}[Threshold localization]
\label{pc:3}
$P\bigl( |\widehat S| \ge r_{NT} \bigr) \to 1$ and
$P\bigl\{ G \ge v_{NT}\Delta_{NT} \ \text{or}\ \widehat S = \emptyset \bigr\} \to 1$,
with $r_{NT}$ as in Assumption~C.3 and $v_{NT}\Delta_{NT} \to 0$,
$v_{NT} \to \infty$.
\end{condition}

Condition~\ref{pc:1} is a statement about \emph{marginal} probabilities under a
fixed reference law: it says the ideal null statistics place little mass in a
band of width $4\Delta_{NT}$ around any admissible threshold, relative to the
tail mass already there.  Condition~\ref{pc:2} says the realized counts track
those marginal probabilities.  Condition~\ref{pc:3} says the data-driven
threshold lands in the range over which the first two conditions are asserted.
Only the combination is needed.

\begin{proposition}[Threshold stability from primitive conditions]
\label{prop:threshold}
Suppose Conditions~\ref{pc:1}--\ref{pc:3} hold and
$(m_{NT} + e_{NT})/r_{NT} \to 0$.  Then Assumption~C.3 holds with
$\gamma_{NT} = C\alpha_{NT}$ and $\varsigma_{NT} \to 0$ for a constant $C$
depending only on the constant in Condition~\ref{pc:2}.
\end{proposition}

\begin{proof}[Proof sketch]
Fix $\epsilon > 0$.  By Condition~\ref{pc:3} it suffices to bound the two
suprema in Assumption~C.3 over $t \in I_{NT}$ on the event
$\{|\widehat S| \ge r_{NT}\}$.

Consider first the regime $M^\star(t) \ge m_{NT}$.  By Condition~\ref{pc:2}
and Chebyshev's inequality applied conditionally on $\mathcal{G}$, both
$M^\star(t)$ and $\widetilde N^\star_{2\Delta}(t)$ concentrate around their
conditional means uniformly over $t \in I_{NT}$: the variance bound divided by
the squared mean is $O(m_{NT}/[\kappa_0\bar G(t)]) + o(\{\ln(NT)\}^{-1})$, and
on $\{M^\star(t) \ge m_{NT}\}$ the first term is $O(1)$ after a standard
chaining over a grid of $t$ of cardinality polynomial in $NT$, with the
$\{\ln (NT)\}^{-1}$ factor absorbing the union bound.  Hence
\[
\sup_{t \in I_{NT},\, M^\star(t)\ge m_{NT}}
\frac{\widetilde N^\star_{2\Delta}(t)}{M^\star(t)}
\;=\;
\sup_{t \in I_{NT}}
\frac{\kappa_0^{-1}\sum_{j} P_{D^\star}(||W_j^\star|-t|\le 2\Delta_{NT}\mid\mathcal{G})}
     {2\bar G(t)}\;\{1 + o_p(1)\}
\;\le\; C\alpha_{NT}
\]
with conditional probability tending to one, where the middle equality uses
$\kappa_0^{-1}\mathbb{E}_{D^\star}[M^\star(t)\mid\mathcal{G}] = 2\bar G(t)$,
which is the null symmetry of $P_{D^\star}$.  This delivers the first bound in
Assumption~C.3 with $\gamma_{NT} = C\alpha_{NT}$.

In the complementary sparse-tail regime $M^\star(t) < m_{NT}$, the band count
is bounded by the same concentration argument by
$e_{NT} := C m_{NT}$ up to $o_p(1)$, and $(m_{NT}+e_{NT})/r_{NT} \to 0$ makes
both negligible relative to the number of discoveries guaranteed by
Condition~\ref{pc:3}.  Collecting the two regimes and the exceptional events
gives $\varsigma_{NT} \to 0$.
\end{proof}

\subsubsection*{Two of the three conditions are implied by existing assumptions}

We now discharge Conditions~\ref{pc:2} and~\ref{pc:3}.

\begin{lemma}[Count concentration from weak score dependence]
\label{lem:pc2}
Work conditionally on $\mathcal{G}$ and on the design-regularity event
$\mathcal{B}_\mathcal{G}$ of Assumption~4.  Let $R = \operatorname{diag}(\Xi)^{-1/2}
\,\Xi\, \operatorname{diag}(\Xi)^{-1/2}$ be the correlation matrix of the
reference score and suppose that, for some $\tau \in (0,1)$ and
$m_{NT}$ as in Condition~\ref{pc:2},
\begin{equation}
\label{eq:weakdep}
\max_{j \le 2\kappa} \bigl| \{ l \le 2\kappa : |R_{jl}| > \tau \} \bigr|
   \;\le\; m_{NT},
\qquad
\max_{j \ne l,\ |R_{jl}| \le \tau} |R_{jl}| \;\le\; C \{\ln (NT)\}^{-1/2}.
\end{equation}
Suppose in addition that on $\mathcal{B}_\mathcal{G}$ the map
$\zeta \mapsto W^\star$ is affine on a neighbourhood of $D + \zeta^\star$ of
radius exceeding $2\Delta_{NT}$, with gradient rows of norm bounded above and
below.  Then Condition~\ref{pc:2} holds.
\end{lemma}

\begin{proof}[Proof sketch]
On the local-affine event, $W_j^\star = a_j' \zeta^\star + b_j$ for
$\mathcal{G}$-measurable $(a_j, b_j)$ with $c \le \|a_j\|_2 \le C$.  Hence
$(W_j^\star)_{j \in \mathcal{H}_0}$ is, conditionally on $\mathcal{G}$, a
Gaussian vector with correlation matrix $\varrho_{jl} = a_j'\Xi a_l /
(\|a_j\|_\Xi \|a_l\|_\Xi)$.  For jointly Gaussian pairs with correlation
$\varrho$, the standard bound
$\bigl|\operatorname{Cov}(\mathbf 1\{W_j^\star \ge t\},
\mathbf 1\{W_l^\star \ge t\})\bigr| \le C|\varrho_{jl}|\,
\bar G_j(t)^{1/2}\bar G_l(t)^{1/2}$ applies.  Splitting the double sum at the
threshold $\tau$ and using~\eqref{eq:weakdep} gives
\[
\operatorname{Var}\Bigl(\sum_{j \in \mathcal{H}_0}\mathbf 1\{W_j^\star \ge t\}\Bigr)
\le \kappa_0 \bar G(t)
 + C m_{NT} \kappa_0 \bar G(t)
 + C \{\ln (NT)\}^{-1/2} [\kappa_0\bar G(t)]^{2}\cdot \kappa_0^{0},
\]
where the second term collects the at most $m_{NT}$ strongly correlated
partners of each $j$ and the third the remaining weakly correlated pairs.
The band-count version is identical with
$\{W_j^\star \ge t\}$ replaced by the band event, whose marginal probability is
bounded by $\bar G(t - 2\Delta_{NT})$.
\end{proof}

\begin{remark}[Why factor adjustment makes~\eqref{eq:weakdep} plausible]
\label{rmk:pca-buys-weakdep}
Condition~\eqref{eq:weakdep} is a weak-dependence requirement on the score
correlation matrix $R$, not on the raw covariates.  It is exactly here that the
factor adjustment of Section~3 pays a second dividend.  The score is built from
the whitened idiosyncratic columns $\breve U_{\widetilde S}$, from which the
common component has been removed by weighted PCA and the double orthogonality
constraint~(3.15).  Under the approximate factor model, the remaining
idiosyncratic cross-correlations are weak by assumption, so $R$ is close to the
Fixed-X pair structure, which is block-diagonal up to the pairwise
$(j, j+\kappa)$ coupling.  Had we constructed knockoffs on the unadjusted
design, the strong cross-sectional correlation induced by the factors would
make~\eqref{eq:weakdep} fail.  Factor adjustment is therefore not only a power
device; it is what makes the threshold-stability route available at all.
\end{remark}

\begin{lemma}[Threshold localization from signal abundance]
\label{lem:pc3}
Suppose Assumptions~1--4 hold, and set $\lambda = C_2 \lambda_{NT}$ as in
Theorem~4.5.  Then Condition~\ref{pc:3} holds with
$r_{NT} = c\kappa$ for the constant $c$ of Assumption~4.
\end{lemma}

\begin{proof}[Proof sketch]
Assumption~4 provides a set $\mathcal{A}_\kappa \subseteq S_R$ with
$|\mathcal{A}_\kappa| \ge c\kappa$ and $|\beta_j| \gg \sqrt{\kappa}\,\lambda_{NT}$
for $j \in \mathcal{A}_\kappa$.  The restricted-eigenvalue and score conditions
of Assumption~4, together with $\chi_{NT} = O_p(\lambda_{NT})$ and
Assumption~8, imply as in the proof of Theorem~4.5 that for every
$j \in \mathcal{A}_\kappa$ the augmented Lasso retains the original coordinate
and kills its knockoff partner with probability tending to one, so that
$W_j^\star \ge c'|\beta_j| \gg \sqrt{\kappa}\,\lambda_{NT}$.  Since
$v_{NT}\Delta_{NT} \to 0$ while $\sqrt{\kappa}\lambda_{NT}$ is bounded below by
the tuning scale, these $c\kappa$ coordinates exceed $v_{NT}\Delta_{NT}$ with
probability tending to one.  The knockoff$+$ ratio in~(3.18) is therefore
already below $q$ at some $g \le \min_{j \in \mathcal{A}_\kappa} W_j^\star$
unless $\widehat S = \emptyset$, which gives both statements.
\end{proof}

\subsubsection*{What remains}

Combining Proposition~\ref{prop:threshold} with Lemmas~\ref{lem:pc2}
and~\ref{lem:pc3}, Assumption~C.3 reduces to Condition~\ref{pc:1} alone.  We
record the form this takes for the LCD statistic and state precisely what is
and is not established.

\begin{remark}[Reduction of Condition~\ref{pc:1} to a conditional density bound]
\label{rmk:density}
On the local-affine event of Lemma~\ref{lem:pc2} the null statistic
$W_j^\star = a_j'\zeta^\star + b_j$ is, conditionally on $\mathcal{G}$,
univariate Gaussian with variance $a_j'\Xi a_j \ge \sigma_{\min}(\Xi)\|a_j\|_2^2
\ge c_0 c^2 > 0$.  Its density is then bounded above by
$(2\pi c_0 c^2)^{-1/2}$, whence
\[
P_{D^\star}\bigl( \bigl||W_j^\star| - t\bigr| \le 2\Delta_{NT} \,\big|\,
\mathcal{G} \bigr) \;\le\; C\,\Delta_{NT}
\qquad \text{uniformly in } t,
\]
and Condition~\ref{pc:1} follows with
$\alpha_{NT} = C \Delta_{NT} / \inf_{t \in I_{NT}} \bar G(t)$, which tends to
zero whenever the null tail at the smallest admissible threshold is of larger
order than the coupling scale.  Since the effective range starts at
$v_{NT}\Delta_{NT}/2$ with $v_{NT}\to\infty$, this is compatible with
$v_{NT}\Delta_{NT}\to 0$.

Two gaps prevent this from being a complete proof, and we do not claim
otherwise.  First, the map $\zeta \mapsto W^\star$ generated by the augmented
Lasso is piecewise affine with kinks at active-set changes; the display above
holds on the event that the active set is locally constant, and a full argument
must control the union over realized active sets rather than fix one.  Second,
the bound degrades near the origin, where a null coordinate may have
$\widehat\beta_j = 0$ while $\widehat\beta_{j+\kappa}$ takes small nonzero
values, so that $W_j^\star$ accumulates just to the right of zero.

The point mass of the LCD statistic at zero, by contrast, is \emph{not} an
obstacle.  The candidate set $\mathcal{W}$ in~(3.18) excludes zero by
construction, and the effective range $I_{NT}$ is bounded away from the origin
at rate $v_{NT}\Delta_{NT}$ with $v_{NT}\to\infty$, so no band in
Assumption~C.3 ever meets the atom.  The residual difficulty is the density in
a shrinking right neighbourhood of zero, not the atom itself.  A complete
treatment of that region, or a switch to a debiased-Lasso coefficient
difference for which the statistic is asymptotically normal with a continuous
limiting density, would close the gap; we leave this to future work.
\end{remark}

\begin{remark}[Relation to the coupling literature]
\label{rmk:ark-comparison}
Threshold conditions of this kind are intrinsic to any coupling-based
robustness argument, because a coupling perturbs a discontinuous selection
rule.  Our Assumption~C.3 is the direct analogue of Conditions~4 and~5 of
Fan, Gao and Lv (2025a): their Condition~5 controls the null mass in a band of
width equal to the coupling accuracy $b_n$ around admissible thresholds,
through the ratio $\{G(t-b_n) - G(t+b_n)\}/G(t)$, and their Condition~4
controls the variance of the null indicator counts.  These correspond
respectively to our Condition~\ref{pc:1} and Condition~\ref{pc:2}, while their
threshold-range lemma corresponds to our Condition~\ref{pc:3}.

The two settings differ in where the randomness that drives the
anti-concentration comes from, and this difference cuts in both directions.
Fan, Gao and Lv verify their Conditions~4 and~5 at a low level for two
statistic constructions, but the verification proceeds through an assumption
on the \emph{feature distribution}: Gaussian covariates in their Section~3.2,
together with sparsity of the precision matrix, which is what controls the
dependence among the null indicators.  They note explicitly that the Gaussian
assumption is imposed mainly to verify these two conditions and can be dropped
if the conditions are assumed directly.  In our conditional-design framework
the reference law $P_{D^\star} = N(D,\Xi)$ is Gaussian \emph{by construction}
rather than by assumption, since we condition on the realized design and work
with the covariance-matched law of the $2\kappa$-dimensional sufficient score;
no distributional restriction on the covariates is required to reach the point
at which an anti-concentration argument can begin.  The counterpart of their
precision-matrix sparsity is our weak-dependence
condition~\eqref{eq:weakdep} on the score correlation matrix, which
Remark~\ref{rmk:pca-buys-weakdep} argues is supplied by the factor adjustment.

Conversely, their coupling operates at the level of the data matrix and is
therefore agnostic to the choice of knockoff statistic, whereas our
score-level reduction (Lemma~4.1) is tied to statistics measurable with respect
to the $2\kappa$-dimensional score.  Their route also benefits from a statistic
with a continuous asymptotic distribution, which is why their verification for
the debiased-Lasso coefficient difference is cleaner than the corresponding
argument would be for the LCD statistic used here.  We regard the two routes as
complementary rather than nested.

One structural feature of our formulation deserves note.  The supremum in
Condition~5 of Fan, Gao and Lv is taken over an interval reaching down to the
origin, whereas Assumption~C.3 restricts attention to $t \ge v_{NT}\Delta_{NT}$
with $v_{NT} \to \infty$.  This truncation, which is legitimate because
Condition~\ref{pc:3} localizes the realized threshold, removes precisely the
region in which a statistic with an atom at zero is least tractable.
\end{remark}
